\chapter{Inventory models}
\label{ch:inventarios}

\section{The decision being modeled}
An inventory policy combines how much to replenish and when to do so. The
minimum cost depends on time units, the demand pattern, and shortage costs;
entering average demand alone is not sufficient
\cite{harris1913many,silver1998inventory}. QSB distinguishes EOQ/EPQ,
quantity discounts, the newsvendor model, period-based lot sizing, and four
stochastic review policies. All use the \texttt{kind}/\texttt{model} JSON
envelope and the native or validation backend. There is no external inventory
backend.

\section{EOQ and gradual production}
With constant demand \(D\), fixed ordering cost \(S\), and holding cost
\(H\) per unit per period, with no shortages, the relevant cost is
\[
 C(Q)=DS/Q+HQ/2,\qquad Q^*=\sqrt{2DS/H}.
\]
Purchasing cost \(Dc\) is added to total cost but does not change the optimum
when unit price \(c\) is constant. Lead time \(\ell\) gives the reorder
point \(r=D\ell\). Rates, lead time, and holding cost must use the same
period. For gradual production at rate \(p>D\), average inventory is
\(Q(1-D/p)/2\), and
\[
 Q^*_{EPQ}=\sqrt{\frac{2DS}{H(1-D/p)}}.
\]
The calculation is continuous: purchasing boxes or whole units requires
comparing feasible quantities. QSB does not automatically impose packaging
constraints.

\subsection{Example: replenishing notebooks}
The file \path{manual/examples/prob_eoq.json} contains \(D=1200\)
units/year, \(S=30\), \(H=2\), \(c=5\), a lead time of \(0.01\)
years, and a current order quantity of \(100\). From the repository root:
\begin{lstlisting}
.build/debug/qsb validate-inventory-json manual/examples/prob_eoq.json
.build/debug/qsb solve-inventory-json manual/examples/prob_eoq.json
\end{lstlisting}
The results are \(Q^*=189.736660\), \(6.324555\) orders/year, a cycle of
\(0.158114\) years, and a reorder point of \(12\). Ordering and holding
each cost \(189.736660\); relevant cost is \(379.473319\), and total
cost is \(6379.473319\). Ordering \(100\) costs \(6460\): the model's
annual saving is \(80.526681\), rather than the entire purchasing cost.
The output document includes the comparison in \texttt{knownQuantity} and
\texttt{closedForm} exactness. In \path{prob_epq.json}, \(p=2400\)
increases the optimal lot to \(268.328157\), with relevant cost
\(268.328157\).

\section{All-units quantity discounts}
QSB evaluates each price tier and its feasible quantities, including purchasing,
setup, and holding costs. The discount applies to every unit in the order; it
is not a marginal block tariff. Holding cost is a cost per unit, rather than a
percentage recalculated from the price.
\path{prob_descuento.json} adds a \(10\%\) discount starting at \(300\)
units. With \texttt{solve-inventory-json}, the selected quantity is
\(Q=300\), with price \(4.5\), setup cost \(120\), holding cost \(300\),
purchasing cost \(5400\), and total cost \(5820\). Relevant cost increases
to \(420\), but the purchasing saving makes this tier preferable. The
candidate table and \texttt{optimum} explain the decision; enumeration of the
tiers is exact within the continuous model.

\section{The newsvendor model}
A single replenishment is chosen for a horizon with no later opportunity to
replenish. For selling price \(v\), purchasing cost \(c\), salvage value
\(s\), and additional shortage penalty \(b\), the marginal costs are
\(C_u=v-c+b\) and \(C_o=c-s\). The critical fractile is
\[
 F(q^*)=\frac{C_u}{C_u+C_o}.
\]
QSB supports normally distributed demand, calculating its quantile and expected
sales, leftovers, shortages, and profit. The probability of covering all demand
is the cycle service level; it is not the expected fraction of demand fulfilled.
In \path{prob_newsboy.json}, mean \(100\), standard deviation \(20\),
\(v=5,c=2,s=1,b=1\), and zero fixed cost give critical fractile \(0.8\),
order quantity \(116.832425\), expected sales \(97.767247\), leftovers
\(19.065178\), shortages \(2.232753\), and profit \(272.003808\).
The normal quantile uses approximate arithmetic; the \texttt{closedForm}
label describes the method, rather than an exact decimal quantity or a guarantee
for nonnormal demand.

\section{Lot sizing}
Ending inventory satisfies \(I_t=I_{t-1}+q_t-d_t\). QSB minimizes
\[
 \sum_t\left(S_t\mathbf{1}_{q_t>0}+c_tq_t+
 h_t\max(I_t,0)+b_t\max(-I_t,0)\right).
\]
Dynamic programming uses integer demand and quantities, zero initial inventory,
and zero final inventory, allowing backorders at a cost. It does not represent
production capacity or configurable initial stock. Exactness refers to the
bounded state space of this educational implementation: inventories between
minus and plus total demand, and production in each period between zero and
that demand. Computational cost grows with total demand and the horizon.
\path{prob_lotes.json} has three demands of ten units, setup cost \(20\),
unit production cost \(2\), holding cost \(1\), and backorder cost \(10\).
QSB returns production \((20,0,10)\), ending inventory \((10,0,0)\), and
cost \(110\). The policy \((10,20,0)\) also costs \(110\); an optimal
result does not imply uniqueness. The decision table allows the balance and
cost allocation to each period to be checked.

\section{Review with uncertain demand}
The four policies use normal demand per period and constant lead time. With
independent demand increments, protection period \(T\) has mean \(DT\)
and standard deviation \(\sigma\sqrt T\). Safety stock is expressed using
the normal quantile, and expected shortage cost uses the normal loss function.
These are continuous-quantity approximations, marked \texttt{approximate};
they do not solve a general optimal policy for arbitrary distributions or
random lead times.

The four examples share annual demand \(1200\), standard deviation \(100\),
lead time \(0.02\) years, setup cost \(30\), holding cost \(2\), backorder
cost \(20\), and purchasing cost \(5\). Every shortage is backordered;
review cost is \(5\). The continuous order-up-to review file uses a supplied
average order size of \(200\).
\begin{center}
\begin{tabular}{lll}
\toprule
File under \texttt{manual/examples/} & Policy & Main result \\
\midrule
\path{prob_qr.json} & Fixed quantity \((Q,r)\) & \(Q=194.866241,r=54.242691\) \\
\path{prob_sS.json} & Continuous order-up-to & \(r=54.095104,S=254.095104\) \\
\path{prob_periodica.json} & Fixed interval & \(T=0.170783,S=321.459989\) \\
\path{prob_opcional.json} & Optional periodic & \(T=0.064550,s=163.951339\) \\
\bottomrule
\end{tabular}
\end{center}
Optional review also returns \(S=353.687998\). Under fixed-quantity continuous
review, safety stock is \(30.242691\), service level \(0.983761\), and
relevant cost \(450.217864\). Under fixed-interval periodic review,
protection includes the review interval plus lead time, \(0.190783\) years:
confusing these intervals artificially reduces the stock required. Backorder
and lost-sales fractions, their costs, and distributions are validated in the
core.

\section{macOS workflow and traceability}
Open the normalized JSON or select an inventory variant in the model creator.
Enter demand, costs, and the variant-specific parameters; validate before
running. Keep units visible and review warnings. In the results, EOQ shows
costs and a comparison with the supplied lot; discounts show candidates; the
newsvendor model shows service and demand outcomes; lot sizing shows quantities
and costs by period; stochastic review shows levels, safety stock, and costs.
The analysis/JSON selector retains the core document for export without
transcribing numbers from charts. Quantities and formulas come from QSBCore
in both the CLI and macOS.

All files in this chapter are validated with \texttt{validate-inventory-json}
and solved with \texttt{solve-inventory-json}; adding
\texttt{--backend validate} validates without solving. The verifier
\path{manual/verify_examples.py} checks the expected results in
\path{manual/examples/manifest.json}; it requires no proprietary reference
material.

\section{Exercises}
\begin{enumerate}
\item Compare integer orders of 189 and 190 in EOQ and explain why comparing
costs is preferable to rounding blindly.
\item Double the lead time in the stochastic example: compare changes in
protected demand and safety stock. Do both double?
\item Lower the backorder cost in the lot-sizing example and check whether
negative inventories appear. Interpret backorders operationally.
\end{enumerate}
